Autonomous mobile robots operating in structured indoor facilities---such as warehouses, hospitals, and office corridors---frequently encounter unmapped or transient physical blockages (such as parked delivery carts or closed safety doors) along nominal routes. Standard navigation systems, including the ROS 2 Navigation Stack (Nav2) \cite{macenski2020marathon2, macenski2022ros2}, utilize layered 2D costmaps \cite{lu2014layered} and local recovery behaviors (e.g., in-place rotations, backup maneuvers, costmap clearing). In standard Nav2 configurations, navigation actions are dispatched independently, and recovery routines operate within the lifetime of an individual goal execution \cite{macenski2020marathon2}. Consequently, when goals are dispatched sequentially from a distant decision junction without cross-goal memory, a purely reactive agent re-enters known dead ends until on-board sensors re-acquire line-of-sight to the blockage.

To mitigate unguided re-exploration, two primary approaches exist:
\begin{enumerate}
    \item \textbf{Spatial Costmap Caching ($O$)}: Maintaining geometric occupancy in persistent global grids (e.g., Costmap2D \cite{lu2014layered}, OctoMap \cite{hornung2013octomap}, Voxblox \cite{oleynikova2017voxblox}) as sensory observations accumulate.
    \item \textbf{Episodic Failure Memory ($F$)}: Recording explicit execution aborts, error causes, or symbolic constraints to suppress infeasible route options in high-level dispatching.
\end{enumerate}

While embodied AI frameworks explore language-based reflection (e.g., Reflexion \cite{shinn2023reflexion}, REFLECT \cite{liu2023reflect}), they are designed for high-level semantic task planning and verbal reflection. In mobile robot navigation, where spatial occupancy is continuously tracked in metric grids, integrating discrete execution failures with sensor-verified invalidation requires explicit coordination with perception. This raises an empirical question: \emph{Under what conditions does an autonomous mobile robot genuinely require explicit episodic failure memory over standard spatial costmap caching?}

In this work, we present \textbf{FailMem}, an event-driven failure memory architecture for ROS 2 Nav2, and conduct an exploratory empirical study to delineate its operational boundaries. Specifically, we provide:
\begin{enumerate}
    \item \textbf{Perception-Grounded Event Memory and Dynamic Invalidation Implementation}: A ROS 2 Nav2 architecture that binds execution aborts to immutable observation bundles (raw LiDAR scans, TF transforms, costmap ROIs) and executes sensor-verified dynamic invalidation upon clearing.
    \item \textbf{Auditable Exploratory Navigation Comparison}: A reproducible 30-run benchmark ($n=3$ per condition under sequential fixed seeds) in Gazebo 11 simulation runs comparing FailMem ($F$) against reactive ($R$), spatial caching ($O$), and persistent suppression ($M1$) baselines across clean (D0), blocked (D1), and restored (D2) scenarios.
    \item \textbf{Boundary Findings on Failure Memory vs. Spatial Caching}: An empirical finding that within the evaluated static 2D hallway layout, specific implementation, and run sample ($n=3$), FailMem yields no observed routing benefit over spatial caching ($O$), selecting identical topological routes ($<1\%$ difference in distance and duration); alongside empirical stopping criteria from an exploratory 4-run feasibility trial where action-conditioned navigation ($H_1$) did not establish executability divergence (Nav2 returned success in 4/4 runs, with strict physical arrival verified in 3/4 runs, without establishing the expected executability divergence).
\end{enumerate}
